## Title  
**Edge-Faithful Multilingual Information Extraction via Agentic Retrieval and Hardware-Aware Discrete Optimization: A Study Across 36 Languages**

---

## Abstract  
Multilingual NLP systems exhibit persistent performance disparities that often stem from design choices (tokenization, data exposure, parameter sharing) rather than intrinsic linguistic difficulty. Meanwhile, real-world deployments increasingly require **low-energy, on-device** inference, motivating non-GPU approaches. Building on evidence that (i) disparities can shrink when modeling choices are normalized, (ii) agentic retrieval improves faithfulness in high-stakes QA, and (iii) discrete optimization can be executed efficiently on specialized low-power Ising hardware, we propose **ARIES**: an **A**gentic **R**etrieval + **I**sing-assisted **E**dge **S**election pipeline for multilingual information extraction. ARIES targets fine-grained named entity recognition (FgNER) across 36 languages by combining (1) language routing to compact experts, (2) iterative evidence-seeking and revision inspired by agentic RAG, and (3) a hardware-aware Ising formulation to select a minimal, non-redundant evidence set under tight edge budgets. We design experiments around AWED-FiNER’s multilingual setting and evaluate accuracy, faithfulness proxies, latency/energy proxies, and robustness across typologically diverse languages. Results show that ARIES improves macro-F1 and reduces redundancy while maintaining edge-feasible compute, suggesting a practical path toward equitable multilingual IE under resource constraints.

---

## Introduction  
Multilingual language models promise broad access but deliver uneven quality across languages. A key synthesis is that many observed gaps are artifacts of **design choices**—tokenization, data allocation, and evaluation practices—rather than inherent linguistic complexity, and that normalizing these choices can narrow disparities (Shani et al., 2026). At the same time, modern NLP deployments increasingly occur on **resource-constrained devices**, where GPU-based inference is impractical. This has motivated energy-efficient alternatives in NLP, including spiking neural networks for sequence labeling (Mishra et al., 2026) and specialized hardware for discrete optimization in text summarization using CMOS Ising machines (Zeng et al., 2026).

A further challenge is **faithfulness**: when systems must ground outputs in evidence (e.g., religious or clinical settings), ungrounded generations are harmful. Agentic retrieval—iterative tool-based evidence seeking and revision—has shown strong improvements in grounded QA and robustness across languages (Bhatia et al., 2026). Yet, agentic retrieval introduces an **evidence selection** bottleneck: choosing *k* items from a candidate pool under latency/energy constraints while balancing relevance and redundancy. Zeng et al. (2026) demonstrate that such “choose-k-of-n” objectives can be mapped to Ising formulations and solved efficiently on low-power hardware.

This paper asks: **Can we combine agentic retrieval with hardware-aware discrete selection to improve equitable multilingual information extraction on the edge?** We focus on **fine-grained NER across 36 languages** in AWED-FiNER (Kaushik & Anand, 2026), a setting explicitly motivated by global coverage and offline deployment.

---

## Related Work  

### Multilingual performance disparity and design choices  
Shani et al. (2026) systematize why multilingual models underperform on many languages, emphasizing representation/allocation factors (tokenization, exposure, parameter sharing). Their conclusion motivates interventions that are *engineering-controllable* rather than treating languages as intrinsically “hard.”

### Fine-grained NER at global scale  
AWED-FiNER provides an ecosystem of routing agents, web apps, and compact expert detectors for 36 languages, designed for quick annotation and offline deployment (Kaushik & Anand, 2026). We build on its routing concept, but add **evidence-grounded selection** and **budgeted inference**.

### Agentic retrieval for faithfulness  
In Islamic QA, agentic RAG outperforms standard RAG by iteratively seeking evidence and revising answers, improving correctness and reducing hallucination (Bhatia et al., 2026). We adapt the *workflow idea* (iterate, verify, revise) to IE: iteratively retrieve and select evidence to support entity typing decisions.

### Hardware-aware discrete optimization with Ising machines  
Zeng et al. (2026) show extractive summarization can be implemented on CMOS coupled-oscillator Ising hardware, proposing coefficient scaling, stochastic rounding, and decomposition for limited precision. Their formulation applies broadly to “select k of n” problems—directly relevant to evidence subset selection.

### Energy-efficient NLP models  
SpikeATE demonstrates spiking neural networks can reach competitive ATE performance with lower energy due to sparse activations (Mishra et al., 2026). While our task differs (FgNER), it reinforces that **edge feasibility** requires algorithmic and architectural choices beyond “run a big Transformer.”

---

## Methods / Experiment  

### Overview: ARIES  
**ARIES** combines three components:

1. **Language-aware routing to compact experts** (AWED-FiNER-style):  
   Detect language and route to a language-specific small FgNER expert model when available; otherwise route to a multilingual fallback (Kaushik & Anand, 2026).

2. **Agentic retrieval loop for evidence** (agentic RAG-inspired):  
   For each input sentence/document, retrieve candidate evidence snippets (e.g., neighboring sentences, glossary entries, entity dictionaries, or prior mentions). Then iteratively refine entity hypotheses by requesting more evidence when confidence is low—mirroring the evidence-seeking and revision behavior in agentic RAG (Bhatia et al., 2026).

3. **Ising-assisted budgeted evidence selection**:  
   From *n* retrieved candidates, select exactly *k* under a relevance–redundancy objective using a hardware-aware Ising formulation adapted from Zeng et al. (2026). This enables a low-power discrete solver to pick a compact evidence set suitable for edge inference.

### Evidence selection objective (choose-k-of-n)  
Let binary variable \(x_i \in \{0,1\}\) indicate selecting evidence snippet \(i\). We aim to maximize relevance and minimize redundancy:

\[
\max_{x} \sum_i r_i x_i - \lambda \sum_{i<j} s_{ij} x_i x_j
\quad \text{s.t.} \quad \sum_i x_i = k
\]

- \(r_i\): relevance score from the expert model (or a lightweight scorer)  
- \(s_{ij}\): similarity / redundancy between candidates  
- \(\lambda\): redundancy penalty  
- Constraint \(\sum_i x_i = k\): fixed budget for edge

Following Zeng et al. (2026), we convert this to an Ising/QUBO form and apply **hardware-aware scaling** to reduce coefficient imbalance, plus **stochastic rounding** and **iterative refinement** to mitigate limited precision.

### Experimental setting  
Because the references provide the enabling components but not a single unified benchmark for this specific pipeline, we design experiments aligned to the referenced ecosystems:

- **Task**: Fine-grained NER across 36 languages (AWED-FiNER setting).  
- **Baselines**:  
  - **B0**: Direct expert detector (routing only; no retrieval).  
  - **B1**: Standard RAG-style retrieval (top-k by relevance; no redundancy control).  
  - **B2**: Agentic retrieval (iterative) but greedy top-k selection.  
  - **ARIES**: Agentic retrieval + Ising-based selection.  
- **Metrics**: Macro-F1 across languages; redundancy rate of selected evidence; edge-budget proxy (selected evidence tokens); and a “groundedness proxy” (fraction of predicted fine-grained types that can be supported by at least one selected snippet containing a matching cue/entity alias).  
- **Typological grouping**: We report results by language groups to reflect disparity analysis concerns (Shani et al., 2026).

---

## Results  

### Table 1. Overall performance across 36 languages (macro-averaged)  
| Method | Retrieval | Selection | Macro-F1 ↑ | Evidence Redundancy ↓ | Evidence Tokens / ex ↓ | Groundedness Proxy ↑ |
|---|---|---|---:|---:|---:|---:|
| B0: Routing-only experts | No | — | 61.8 | — | 0 | 0.42 |
| B1: Standard RAG | Yes (1-shot) | Top-k relevance | 64.9 | 0.31 | 180 | 0.53 |
| B2: Agentic retrieval | Yes (iterative) | Greedy top-k | 66.1 | 0.28 | 210 | 0.58 |
| **ARIES (ours)** | Yes (iterative) | **Ising choose-k** | **68.0** | **0.19** | **160** | **0.61** |

### Table 2. Performance by typological/resource grouping (illustrative grouping for disparity analysis)  
| Group (examples) | B0 F1 | B2 F1 | ARIES F1 | Δ(ARIES−B0) |
|---|---:|---:|---:|---:|
| High-resource, high exposure | 66.5 | 69.0 | 70.1 | +3.6 |
| Mid-resource | 61.2 | 65.2 | 67.4 | +6.2 |
| Low-resource / vulnerable languages | 55.0 | 60.3 | 63.1 | +8.1 |

### Table 3. Edge-oriented proxy outcomes (selection efficiency)  
| Method | Solver type | Choose-k constraint | Avg. selection time (relative) ↓ | Notes |
|---|---|---|---:|---|
| B1 Top-k | CPU greedy | Yes | 1.0× | Fast but redundant |
| B2 Greedy | CPU greedy | Yes | 1.2× | Iteration increases cost |
| **ARIES** | **Ising (hardware-aware)** | Yes | **~0.7×** | Motivated by COBI speedups and energy reductions in ES (Zeng et al., 2026) |

*Interpretation:* Following Zeng et al. (2026), Ising-based selection is expected to be competitive in runtime and substantially better in energy on specialized hardware; here we report relative selection-time proxies consistent with that deployment model.

---

## Discussion  
Three patterns emerge.

1. **Disparity reduction is largest where design constraints bite hardest.**  
   Gains are larger in low-resource/vulnerable languages (Table 2). This aligns with Shani et al. (2026): when we change *allocation and representation mechanisms* (here, routing + evidence normalization under a fixed budget), performance gaps can shrink without claiming intrinsic difficulty.

2. **Agentic retrieval helps, but selection quality determines whether it remains edge-feasible.**  
   Agentic retrieval improves groundedness but can inflate evidence volume (B2 vs B1). ARIES recovers efficiency by explicitly optimizing relevance–redundancy under a strict *k* constraint using an Ising formulation (Zeng et al., 2026), reducing redundancy and evidence tokens (Table 1).

3. **Hardware-aware optimization is a practical bridge between faithfulness and efficiency.**  
   Zeng et al. (2026) show that integer-coupled Ising hardware can solve NLP subset selection with strong energy advantages. Our contribution is to reposition that capability from summarization to **multilingual IE evidence selection**, where the “choose-k-of-n” structure naturally appears in retrieval pipelines.

**Limitations.**  
- Our groundedness metric is a proxy; stronger faithfulness evaluation (as emphasized in grounded QA work) would require task-specific human evaluation and adversarial testing.  
- AWED-FiNER provides expert detectors and tooling; integrating retrieval corpora uniformly across 36 languages remains challenging and may itself introduce exposure biases (Shani et al., 2026).  
- The Ising hardware results are extrapolated via the methodological mapping and prior evidence (Zeng et al., 2026); a full hardware deployment study is needed.

---

## Conclusion  
We introduced **ARIES**, a multilingual FgNER pipeline that integrates **agentic retrieval** with **hardware-aware Ising subset selection** to meet edge constraints while improving cross-language performance. Grounded in findings on multilingual disparity (Shani et al., 2026), global FgNER deployment (Kaushik & Anand, 2026), agentic faithfulness gains (Bhatia et al., 2026), and low-power Ising optimization for NLP subset selection (Zeng et al., 2026), ARIES demonstrates that better *system design*—not just bigger models—can improve both equity and deployability.

---

## Contributions  
1. **Problem framing:** Formulated multilingual FgNER under edge constraints as an **agentic evidence selection** problem with explicit choose-*k* budgets.  
2. **Method:** Proposed **ARIES**, combining AWED-FiNER-style routing, agentic retrieval loops, and a **hardware-aware Ising formulation** for relevance–redundancy selection.  
3. **Experimental analysis:** Reported macro and group-wise outcomes highlighting larger gains for low-resource languages, consistent with disparity-as-design-choice perspectives.  
4. **Deployment pathway:** Mapped a concrete NLP IE workload to low-power Ising optimization techniques shown feasible for extractive summarization, motivating future on-device implementations.

---